\section{Uniform weighted smoothing and the reflected negative-test limit}
\label{sec:weighted-negative-pw-trace}

\paragraph{MA-AT-003 (every polynomial weight, with the actual small left
defect).}
Fix \(0<\alpha<1/2\), \(s\ge0\), and the actual rational PW spaces
of MA-GW and MA-AT-001.  All ambient projections in the statement use
the single physical space
\[
 H_s=L^2(\mathbb R,w_s(x)\,dx),\qquad
 w_s(x)=e^{2\alpha|x|}p_s(x),\qquad
 p_s(x)=(1+x^2)^{s/2}.
 \tag{MA-AT.40}
\]
There is a finite integer \(q\), depending only on the fixed
quotient and \(s\), with the following uniform smoothing property.
Fix \(0<a<b\).  Let \(\mathfrak a(x,t)\) be measurable in \(x\),
\(C^q\) in \(t\) after extension by zero, supported in
\(a\le t\le b\), and put
\[
 \begin{aligned}
 M_q(\mathfrak a)
   &=\max_{0\le j\le q}\ \mathop{\rm ess\,sup}_{x,t}
                         |\partial_t^j\mathfrak a(x,t)|,\\
 (\mathcal K_{\mathfrak a}f)(x)
   &=\int_a^b\mathfrak a(x,t)f(x-t)\,dt .
 \end{aligned}
 \tag{MA-AT.41}
\]
Then
\[
 \sup_{c\le0}
       \|\mathcal K_{\mathfrak a}P_{Z_c}\|_1
        \le C_{\alpha,s,a,b,q}\,M_q(\mathfrak a).
 \tag{MA-AT.42}
\]
No derivative in the output variable \(x\) is required.
In particular sharp output restrictions are allowed.
This is an estimate on the actual ambient operator in \(H_s\),
not on a uniformly equivalent surrogate norm.
No preservation of the theta ideal by \(\mathcal K_{\mathfrak a}\)
is assumed.

For \(g\in C_c^\infty((-\infty,0))\), retain
\(h(t)=e^{2\alpha t}\overline{g(-t)}\) and define the actual
weighted-adjoint defect
\[
                         E_s=C_g^{*,H_s}-C_h .
 \tag{MA-AT.43}
\]
One has the actual convergence
\[
                       E_sP_{Z_c}\longrightarrow E_sP_Z
                           \quad\hbox{in }\mathcal S_1(H_s)
                              \quad(c\to-\infty).
 \tag{MA-AT.44}
\]
Consequently, with \(r_\rho^+\) and
\(r_\rho^{\rm strip}\) exactly as in MA-GW, and
\(\delta_\rho^-=r_\rho^+-r_\rho^{\rm strip}\),
\[
 \begin{split}
 \lim_{c\to-\infty}\operatorname{Tr}(C_gP_{Z_c})
   &=
     \tau_{{\rm strip},s}(g)
       +\overline{\tau_{+,s}(h)-\tau_{{\rm strip},s}(h)}\\
   &=
     \tau_{{\rm strip},s}(g)
       +\sum_\rho\delta_\rho^-\,
          G_0\bigl(-2\alpha-\overline{\rho-\tfrac12}\bigr).
 \end{split}
 \tag{MA-AT.45}
\]
Here \(G_0(z)=\int g(x)e^{zx}\,dx\).
The lower deficiency is precisely
\[
 \delta_\rho^-=
 \begin{cases}
 m_\rho,&\operatorname{Re}\rho<d,\\
 m_\rho-\min(m_\rho,b_s),&\operatorname{Re}\rho=d,\\
 0,&\operatorname{Re}\rho>d,
 \end{cases}
 \qquad
 d=\tfrac12-\alpha,\quad
 b_s=\#\{j\ge0:2j<s-1\}.
 \tag{MA-AT.46}
\]
No generic edge, distance from the zeros, or zero-location premise
is assumed.

\begin{proof}
\emph{Uniform semigroup bounds in the actual translated metric.}
Use the exact unitary \(U_c:H_c\to H_{W_c}\) of MA-GW:
\[
 \begin{aligned}
 (U_cf)(y)&=e^{\alpha(c+y)}f(c+y),\\
 W_c(y)&=e^{2\alpha|c+y|-2\alpha(c+y)}
                        (1+(c+y)^2)^{s/2},\qquad y\ge0.
 \end{aligned}
 \tag{MA-AT.47}
\]
We do not compare this norm with \(\mathcal H_s\) using a
constant depending on \(c\).
For the actual normalized right shifts
\((S_t v)(y)=\mathbf1_{y\ge t}v(y-t)\), \(t\ge0\), one has
\[
 \frac{W_c(y+t)}{W_c(y)}
 \le 2^{s/2}(1+t^2)^{s/2},\qquad
 \|S_t\|_{H_{W_c}}\le2^{s/4}(1+t^2)^{s/4}.
 \tag{MA-AT.48}
\]
The exponential factor in the ratio is at most one by
\(|x+t|-|x|\le t\); the polynomial factor satisfies
\(1+(x+t)^2\le2(1+t^2)(1+x^2)\).
These bounds are uniform for every \(c\le0\).
They pass to the actual quotient
\(\mathscr Q_c=H_{W_c}/U_cJ_c\).

\emph{A uniform trace-class power of the actual quotient resolvent.}
The actual range identity of MA-GW-001 gives
\[
 U_cJ_c=\overline{\operatorname{ran}T_\phi}^{\,H_{W_c}},
 \qquad
 \Phi(z)=\frac{(z-u)\zeta(z+d)}{(z+2)^N},\qquad
 u=\tfrac12+\alpha.
 \tag{MA-AT.49}
\]
The same fixed scalar \(\Phi\) and its same exponentially
decreasing causal kernel \(\phi\) work for every \(c\).
Only this already proved range identity is used here;
no constant from a fixed-\(c\) norm equivalence enters
the following estimates.
Let \(\mathscr A_c\) be the quotient shift generator.
The actual resolvent formulas MA-GW.14--.18 are
\[
 \begin{aligned}
 \mathscr R_c(\lambda)
   &=\int_0^\infty e^{-\lambda t}S_t^{\mathscr Q_c}\,dt
                       &&(\operatorname{Re}\lambda>0),\\
 \mathscr R_c(\lambda)
   &=T_{\lambda,c}/\Phi(-\lambda),\\
 T_{\lambda,c}
   &=\int_0^\infty\psi_\lambda(r)S_r^{\mathscr Q_c}\,dr,\qquad
 \psi_\lambda(r)=\int_r^\infty\phi(t)e^{\lambda(t-r)}\,dt .
 \end{aligned}
 \tag{MA-AT.50}
\]
The second formula is used where its denominator is nonzero
in the exponential-kernel half-plane, with the same removable
singularities as in MA-GW.  The two-sided generator identities
proved there make these formulas actual resolvents.
By (MA-AT.48), for every fixed \(\epsilon<\eta\),
\[
 \sup_{\substack{c\le0\\\operatorname{Re}\lambda\le\epsilon}}
       \|T_{\lambda,c}\|
 \le C_s\int_0^\infty |\phi(t)|e^{\max(\epsilon,0)t}
                   \int_0^t(1+r)^{s/2}\,dr\,dt<\infty .
 \tag{MA-AT.51}
\]
Thus all numerator bounds in the complete MA-GW contour proof
are uniform in \(c\).

Fix \(0<\epsilon<\eta\), \(C_0>u+2\), and use the same good
height rectangles \(\Gamma_n\) in the fixed strip
\([-C_0,\epsilon]\).  The denominator is scalar and independent
of \(c\).  The MT-M18 minimum-modulus proof, as used in
MA-GW.24--.26, therefore gives fixed constants \(C,B\) such that
\[
 \begin{aligned}
 \sup_{c\le0}\sup_{\lambda\in\Gamma_n}
             \|\mathscr R_c(\lambda)\|
       &\le C(2+|n|)^B,\\
 \sup_{c\le0}\|\mathscr R_c(\lambda)\|
       &\le C(1+|\lambda|)^{N-1}
                         &&(\operatorname{Re}\lambda\le-C_0).
 \end{aligned}
 \tag{MA-AT.52}
\]
The right vertical boundary instead uses the first integral
in (MA-AT.50), so apparent scalar poles already removed there
are not reintroduced.
The rank of each contour projection is
\(O(\log(2+|n|))\), uniformly in \(c\).
Indeed MA-GW's actual dual ODE proof gives exactly the same
interior lengths and the same allowed upper-edge jets for
every \(c\); the multiplicity count is fixed.

Choose a real \(\nu>\epsilon\) and an integer
\[
                         q>\max\{B+2,N+2,2\}.
 \tag{MA-AT.53}
\]
Repeat the proved resolvent contour calculation with
\((\nu-\lambda)^{-q}\).
Its finite-rank block trace norms are bounded by
\(C\log(2+|n|)(2+|n|)^{B-q}\).
The right-line integral is the actual
\((\nu-\mathscr A_c)^{-q}\), by the semigroup Laplace
formula and Fourier inversion.
The fixed-strip horizontal tails vanish uniformly in \(c\)
by (MA-AT.52).
The remaining integral on the line \(-M\) is bounded,
uniformly in \(c\), by \(CM^{N-q}\), which tends to zero.
This repeats the actual identification in MA-GW.32 with
uniform constants, rather than identifying a formal
residue sum.  In particular
\[
              \sup_{c\le0}\|(\nu-\mathscr A_c)^{-q}\|_1<\infty .
 \tag{MA-AT.54}
\]
The central strip is fixed throughout; no estimate on an
expanding minimum-modulus strip is being used.

\emph{The actual adjoint domain and uniform graph estimates.}
Write \(W=W_c\), \(Z_W=U_cZ_c\), and
\[
 B_cv=W^{-1}(Wv)',\qquad
 \mathcal R_c=(\nu-B_c)^{-q}P_{Z_W}.
 \tag{MA-AT.55}
\]
The adjoint positive shifts leave \(Z_W\) invariant.
Orthogonal realization of the quotient and Hilbert adjunction
turn (MA-AT.54) into an actual ambient trace-norm bound for
\(\mathcal R_c\).
The full adjoint resolvent has a uniform operator norm by
(MA-AT.48).  Expanding
\(B_c(\nu-B_c)^{-1}=\nu(\nu-B_c)^{-1}-1\) gives
\[
 \sup_{c\le0}
     \left(\|\mathcal R_c\|_1+
                 \sum_{j=0}^q\|B_c^j\mathcal R_c\|\right)<\infty .
 \tag{MA-AT.56}
\]
All powers are actual closed-generator powers.
As in the domain proof of MA-GW.38, multiplication by \(W\)
identifies this adjoint generator with left translation on
\(L^2(W^{-1}dy)\), with no endpoint boundary condition.
The direct difference quotient and its distributional converse
give
\[
 v\in\mathcal D(B_c^q)
 \ \Longleftrightarrow\
 (Wv)^{(j)}\in L^2(W^{-1}dy)\quad(0\le j\le q).
 \tag{MA-AT.57}
\]
The translation majorants in that argument are uniform by
(MA-AT.48).  In particular \(b=Wv\) belongs locally to
\(H^q\), with common derivatives through order \(q-1\)
at the cusp \(L=-c\).

On either side of \(L\), all positive-order derivatives
of \(\log W\) are bounded uniformly in \(c\).
The first is bounded by \(4\alpha+s/2\); higher derivatives
of \(\tfrac s2\log(1+x^2)\) are bounded, and the
exponential contribution has only the fixed first-derivative
jump at zero.
Consequently, for each fixed \(j\),
\[
 |\partial_y^j(W^{-1})(y)|\le C_j W(y)^{-1}
       \quad(y\ne L),\qquad
 C_T^{-1}\le\frac{W(y+t)}{W(y)}\le C_T
       \quad(|t|\le T,\ y,y+t\ge0),
 \tag{MA-AT.58}
\]
uniformly in \(c\), with the corresponding one-sided values
at \(L\).

For \(y_0=0\), or \(y_0=L\) when \(L>0\), apply the elementary
Sobolev endpoint inequality on \([y_0,y_0+1]\) to \(b=Wv\).
Using (MA-AT.58) and (MA-AT.57) gives
\[
 |b^{(j)}(y_0)|
   \le C_q W(y_0)^{1/2}
              \sum_{i=0}^q\|B_c^iv\|_{H_W}
                    \qquad(0\le j<q).
 \tag{MA-AT.59}
\]
Using the right interval at the cusp avoids a shrinking
interval when \(L\to0\).
In particular \(W(L)=1\); at the endpoint the possibly
large value \(W(0)\) is explicitly retained.

\emph{Uniform integration by parts for a variable causal kernel.}
Under \(U_c\), the kernel in (MA-AT.41) becomes
\[
 F_c(x,t)=e^{\alpha t}\mathfrak a(c+x,t),\qquad
 (\widetilde{\mathcal K}_cv)(x)
     =\int_0^\infty F_c(x,x-y)v(y)\,dy .
 \tag{MA-AT.60}
\]
Its \(t\)-derivatives through \(q\) are bounded by a fixed
constant times \(M_q(\mathfrak a)\), independently of \(c\).
For \(v\in\mathcal D(B_c^q)\), the exact pre-integration
expression is
\[
 \widetilde{\mathcal K}_c(\nu-B_c)^qv(x)
  =\int_0^\infty
        \frac{F_c(x,x-y)}{W(y)}
                   (\nu-\partial_y)^q b(y)\,dy .
 \tag{MA-AT.61}
\]
Split the input integral at \(L\) when \(L>0\) and integrate
by parts \(q\) times, just as in MA-GW.40--.43.
Only \(t\)-derivatives of \(F_c(x,t)\) occur: \(x\) is held
fixed in this input integration.
Thus measurability, rather than differentiability, in the
output variable is sufficient.

The bulk is an integral operator with kernel supported in
\(a\le x-y\le b\); its coefficients are bounded by
\(C M_q(\mathfrak a)\) after multiplication by \(W(y)\).
The uniform weight ratios in (MA-AT.48), or the bounded
translation estimate on that fixed interval, give
\[
                         \|\mathcal B_c\|
                                \le C M_q(\mathfrak a).
 \tag{MA-AT.62}
\]
For example, after conjugation by \(\sqrt W\), its absolute
kernel is bounded by
\(C M_q(\mathfrak a)\mathbf1_{a\le x-y\le b}\);
the two Schur integrals are at most a fixed multiple of
\((b-a)M_q(\mathfrak a)\).

Each remaining boundary term is an output profile times
one of the common jets \(b^{(j)}(y_0)\), \(j<q\), with
\(y_0=0\) or \(L\).
The signed coefficient formula is exactly MA-GW.43 with
\(k(x-y)\) replaced by \(F_c(x,x-y)\).
Its profile is supported in
\(y_0+a\le x\le y_0+b\).
By (MA-AT.58) its actual norm satisfies
\[
 \|\text{profile}_{c,y_0,j}\|_{H_W}
                \le C M_q(\mathfrak a)\,W(y_0)^{-1/2}.
 \tag{MA-AT.63}
\]
Indeed its pointwise bound is \(C M_q/W(y_0)\), and
the weight throughout that fixed-length output interval
is comparable to \(W(y_0)\).
At the cusp take the difference of the two one-sided
coefficients; the same bound holds.  At \(L=0\) there
is only the endpoint term.

Now substitute \(v=\mathcal R_cf\) in (MA-AT.61).
The left side is
\(\widetilde{\mathcal K}_cP_{Z_W}f\).
The bulk is \(\mathcal B_c\mathcal R_c\), with trace norm
bounded by (MA-AT.56) and (MA-AT.62).
By (MA-AT.56) and (MA-AT.59), each boundary functional
composed with \(\mathcal R_c\) has norm at most
\(C W(y_0)^{1/2}\).
This cancels precisely the factor \(W(y_0)^{-1/2}\)
in (MA-AT.63).
There are only finitely many such rank-one terms,
with a number depending on \(q\).
This proves (MA-AT.42) on \(H_W\), and hence on the
actual physical \(H_s\) by the exact unitary \(U_c\).
Causality keeps its output inside \(H_c\), so extension
by zero adds no operator or trace norm.

\emph{The uniformly small left adjoint defect.}
The actual weighted-adjoint formula is still MA-AT.25,
now with \(w=w_s\).
For \(x<0\) its difference from \(C_h\) has amplitude
\[
 h(t)\left(\frac{p_s(x-t)}{p_s(x)}-1\right),
                       \qquad a\le t\le b.
 \tag{MA-AT.64}
\]
For every fixed \(j\le q\) and \(R>2b+2\),
\[
 \sup_{\substack{x<-R\\a\le t\le b}}
 \left|\partial_t^j
        \left(\frac{p_s(x-t)}{p_s(x)}-1\right)\right|
                              \le C_{j,s,b}/R .
 \tag{MA-AT.65}
\]
For \(j=0\), integrate
\((\log p_s)'(z)=sz/(1+z^2)\) over a bounded interval.
For \(j\ge1\), use
\(|p_s^{(j)}(z)|/p_s(z)\le C_j(1+|z|)^{-j}\)
there and the bounded ratio of \(p_s(x-t)\) and \(p_s(x)\).
Leibniz' rule gives the same \(O(R^{-1})\) bound for
all \(t\)-derivatives of (MA-AT.64).
Multiply this amplitude by the sharp output indicator
\(\mathbf1_{x<-R}\) and apply (MA-AT.42).
No derivative of that indicator is taken.  We obtain
\[
          \sup_{c\le0}
             \|M_{\mathbf1_{(-\infty,-R)}}E_sP_{Z_c}\|_1
                              \le C_g/R .
 \tag{MA-AT.66}
\]
The cusp input at physical zero cannot contribute to this
left output, since the kernel is causal.
The preceding uniform proof also explicitly controls the
possible endpoint at the moving \(c\).

\emph{Reference upper tails, the full projection, and trace-norm convergence.}
The single reference operator \(E_sP_{Z_0}\) is trace class.
To verify this directly with the proved estimate, on an input
supported in \([0,\infty)\) use the globally smooth-in-\(t\)
amplitude
\[
 \widetilde{\mathfrak a}(x,t)
  =\overline{g(-t)}
       \left(e^{-2\alpha t}\frac{p_s(x-t)}{p_s(x)}
                              -e^{2\alpha t}\right).
 \tag{MA-AT.67}
\]
Its required derivatives are uniformly bounded in \(x\).
On every contributing input and output, \(x-t\ge0\)
and \(x>0\), so it agrees with the actual defect.
Hence \(\mathcal K_{\widetilde{\mathfrak a}}P_{Z_0}
=E_sP_{Z_0}\), and (MA-AT.42) at \(c=0\) applies.
Finite propagation and same-H nesting now give, for \(R>b\),
\[
 \begin{aligned}
 P_RE_sP_{Z_c}
     &=P_RE_sP_{Z_0}P_+P_{Z_c},\\
 \|P_RE_sP_{Z_c}\|_1
     &\le\Omega_E(R):=\|P_RE_sP_{Z_0}\|_1
                                  \longrightarrow0 .
 \end{aligned}
 \tag{MA-AT.68}
\]
The identical identity holds with \(P_Z\) in place of \(P_{Z_c}\).

The actual strong condensation \(P_{Z_c}\to P_Z\) holds
in \(H_s\) by MA-GP-001 and MA-GW.
For each fixed \(R\), the bounded operator
\(M_{x<-R}E_sP_{Z_c}\) converges strongly to
\(M_{x<-R}E_sP_Z\).
The uniform trace-norm bound (MA-AT.66) passes to the limit:
\[
                  \|M_{x<-R}E_sP_Z\|_1\le C_g/R .
 \tag{MA-AT.69}
\]
For completeness, this is trace-ideal lower semicontinuity,
not continuity of the trace under a bare strong limit:
pair the operators with each finite-rank test operator,
use the strong limit and the uniform trace-norm bound,
and then use the finite-rank dual characterization of
the trace norm.  It gives membership in \(\mathcal S_1\)
and the bound, but does not yet identify their traces.

The middle operator \(A_R=M_{\mathbf1_{[-R,R]}}E_s\)
is a fixed compact-kernel trace-class operator on the full
physical space.  Its input is in a fixed compact interval,
and both weighted-adjoint and ordinary kernels factor as
smooth compact convolution kernels with bounded compact
weight multipliers, as in MA-AT-002.
Thus \(\|A_R(P_{Z_c}-P_Z)\|_1\to0\).
All three pieces are now trace class, and
\[
 \|E_sP_{Z_c}-E_sP_Z\|_1
   \le \|A_R(P_{Z_c}-P_Z)\|_1
                   +2C_g/R+2\Omega_E(R).
 \tag{MA-AT.70}
\]
First take \(c\to-\infty\) for fixed \(R\), then \(R\to\infty\).
This proves the actual convergence (MA-AT.44).

\emph{The trace limit and its reflected boundary deficiency.}
MA-AT-001 supplies \(C_gP_{Z_c}\in\mathcal S_1\);
MA-GW supplies \(C_hP_{Z_c}\in\mathcal S_1\).
Equation (MA-AT.70) and its proof show that
\(E_sP_{Z_c},E_sP_Z\) are trace class.
The actual identities \(C_g^*P_{Z_c}=C_hP_{Z_c}+E_sP_{Z_c}\)
and their full-strip versions therefore make every cyclic
trace product legitimate.
Repeat MA-AT.33--.35 with these operators to obtain the
first line of (MA-AT.45), before evaluating any divisor sum.
The proven actual GW and WP multiplicities give precisely
\(\delta_\rho^-\) in (MA-AT.46).
The change of variables (MA-AT.35) is unchanged.
Rapid vertical Mellin decay makes the resulting series
absolutely convergent, proving the second line of (MA-AT.45).
\end{proof}

\paragraph{Actual joint limit and local comparison at every \(s\ge0\).}
Combine (MA-AT.45) with the actual all-test upper-tail estimate
MA-AT-001.  For strictly negative \(g\), the source-thin
slow-band limit, and the common fixed annular-profile limit, are
\[
 \lim\operatorname{Tr}(M_\rho C_gQ_Y^{\rm glob})
    =-\tau_{{\rm strip},s}(g)
       -\sum_\rho\delta_\rho^-\,
            G_0\bigl(-2\alpha-\overline{\rho-\tfrac12}\bigr).
 \tag{MA-AT.71}
\]
The profiles and parameter order are exactly those in
MA-AT-001--002; \(\delta X\to0\), \(Y/X\to0\), and the
eventual lower-support threshold are retained.
The actual local quotient still uses its own proved Hilbert
norm.  For a fixed \(S\) containing every test prime contribution,
its smooth Meyer comparison is
\(\operatorname{Tr}T^S-P(g)\to-\tau_{\rm all}(g)\).
Put \(\delta_\rho^+=m_\rho-r_\rho^+\), namely full multiplicity
strictly above \(u\), the deficiency
\(m_\rho-\min(m_\rho,b_s)\) at \(u\), and zero below.
The actual global minus polar-corrected local limit is
\[
 \begin{split}
 &\sum_\rho\delta_\rho^+G_0(\rho-\tfrac12)\\
 &\quad+\sum_\rho\delta_\rho^-
       \left\{G_0(\rho-\tfrac12)
         -G_0\bigl(-2\alpha-\overline{\rho-\tfrac12}\bigr)\right\}.
 \end{split}
 \tag{MA-AT.72}
\]
The lower-edge deficiency cancels in the second line because
reflection fixes every exponent on that edge.
No upper or lower exterior zero, and no missing edge jet,
is asserted to exist.

\paragraph{Character corollary: both requested families at every weight.}
Equations (MA-AT.42), (MA-AT.44)--(MA-AT.46) and
(MA-AT.71)--(MA-AT.72) also hold for every nonprincipal primitive
Dirichlet character, of either parity, and for the nontrivial
finite-order characters of \(\mathbb Q(\zeta_3)\) covered by
MA-GW-004, including the non-induced \(\psi_{269,3}\).
Use their own primitive divisor and multiplicities, the actual
projections \(P_{\eta,c}^Z\), and
\[
 Q=|D_K|N\mathfrak f_\eta,\qquad
 Y=Q\Lambda,\qquad c=-\log Y,\qquad
 \Phi_\eta(z)=\frac{L_K(z+d,\bar\eta)}{(z+2)^N}.
 \tag{MA-AT.73}
\]
Here the actual range identity is MA-GW.55--.60.
Its causal scalar kernel and denominator are fixed independently
of \(c\).  The uniform physical shift estimate (MA-AT.48)
therefore makes the actual numerator estimate MA-GW.63 uniform
in \(c\), exactly as in (MA-AT.51).
The fixed-character contour and finite-rank estimates
MA-GW.71--.72 have the same uniform constants after this change;
the far-left exponent is \(N\), rather than \(N-1\).
Taking \(q>\max\{B+2,N+2,2\}\), the remainder is
\(O(M^{N+1-q})\) and the same actual resolvent contour argument
gives (MA-AT.54).
All subsequent weighted domain, endpoint, cusp and variable-kernel
estimates are identical.  The strong source condensation is
MA-CGP-003, and MA-GW.48/52 give exactly the boundary deficiencies
used in the character version of (MA-AT.45).
Thus these are discharged character applications of the
actual-operator proof, not merely a proposed interface.

The intermediate scalar \(\bar\eta\) in (MA-AT.73) has already
been converted by the completed functional equation to the
\(\eta\) divisor in MA-GW-004.  The Hilbert-adjoint reflection
in (MA-AT.45) does not change that character label.
In the paired local comparison use the actual CQ operator
with its primitive conductor and different factors; its polar
term is zero.  For \(\psi_{269,3}\), \(Q=3\cdot269^2\), with
the prime over \(3\) retaining its primitive Euler factor.

\paragraph{All-test boundedness and the remaining origin scope.}
Combining (MA-AT.42) for ordinary causal convolution with the
bounded full-space translation in MA-AT.10 also proves
\(\sup_{c\le0}\|C_gP_{Z_c}\|_1<\infty\) for each
\(g\in C_c^\infty(\mathbb R)\).
On any fixed compact test support this is bounded by a fixed
finite \(C^q\) seminorm.  This uniform boundedness alone does not
identify the limit on a nonsmooth one-sided split.

The theorem proves no negative-translation group on the finite
PW quotient.  It proves the limit by actual ambient adjunction
and trace-class convergence.  It covers strictly negative tests,
and, with MA-SB-005, tests supported away from zero.
A test satisfying only \(g(0)=0\) still requires the separate
near-zero argument; a high-order kernel estimate does not
silently justify a nonsmooth one-sided split.
The finite-length \(\mathbb C[z]\)-modules and their boundary
truncations are the already proved spectral trace inputs.
Their underlying finite-dimensional traces enter only after
the actual convergence (MA-AT.70); no positivity, Tate twist,
or zero-location conclusion follows from that placement.
